%% kapa: exam=1 round=35 year=2024 period=1 form=A subject=civil_law subject_ko=민법 seq=25 number=25
%% source: Q-Net 감정평가사 기출문제 (cst00309) artlSeq=5213229 file=raw/1008/5213229/1교시 문제지 원본.pdf
%% conversion: kordoc 4.24.51 -> tools/gen1.py (auto); review=proofread-sonnet; flags=none
\begin{question}{25}
점유자와 회복자의 관계에 관한 설명으로 옳은 것은? (다툼이 있으면 판례에 따름)
\begin{choices}
  \item 지상권자는 선의점유자라도 자주점유자가 아니므로 과실수취권이 인정되지 아니한다.
  \item 타주점유자가 점유물을 반환하는 경우, 점유자는 특별한 사정이 없는 한 회복자에 대하여 점유물을 보존하기 위하여 지출한 금액의 상환을 청구할 수 있다.
  \item 악의의 점유자는 과실(過失)없이 과실(果實)을 수취하지 못한 경우에도 그 대가를 보상하여야 한다.
  \item 점유물이 점유자의 책임있는 사유로 멸실된 경우, 선의의 타주점유자는 이익이 현존하는 한도에서 배상하여야 한다.
  \item 점유자가 점유물에 유익비를 지출한 경우, 특별한 사정이 없는 한 점유자는 회복자에 대하여 그 가액의 증가가 현존한 경우에 한하여 점유자의 선택에 좇아 그 지출금액이나 증가액의 상환을 청구할 수 있다.
\end{choices}
\end{question}
\answer{2}
